% Part: first-order-logic
% Chapter: natural-deduction
% Section: proving-things

\documentclass[../../../include/open-logic-section]{subfiles}

\begin{document}

\iftag{FOL}
      {\olfileid{fol}{ntd}{pro}}
      {\olfileid{pl}{ntd}{pro}}

\olsection{\usetoken{P}{derivation}కు ఉదాహరణలు}

\begin{ex}
\tetoken{వాక్యం}{!!{sentence}}~$(!A \land !B) \lif !A$ యొక్క ఒక \tetoken{వ్యుత్పత్తిని}{!!a{derivation}}
ఇద్దాం.

కావలసిన నిష్కర్షను \tetoken{వ్యుత్పత్తి}{!!{derivation}} అడుగున రాయడంతో మొదలుపెడతాం.
\begin{prooftree}
\AxiomC{}
\UnaryInfC{$(!A\land !B) \lif !A$}
\end{prooftree}

ఇలాంటి రూపంలోని ఒక \tetoken{వాక్యాన్ని}{!!a{sentence}} ఏ రకపు నిగమనం ఇవ్వగలదో ఇప్పుడు
కనుగొనాలి. నిష్కర్షలోని \tetoken{ప్రధాన సంయోజకం}{!!{main operator}}~$\lif$; కనుక
\Intro{\lif} నియమంతో నిష్కర్షకు చేరడానికి ప్రయత్నిద్దాం. ముందుకు
సాగుతున్నప్పుడే సంబంధిత పరికల్పనలను రాసి, నిగమన నియమాలకు గుర్తులు
పెట్టడం మంచిది. అప్పుడు నిరూపణ చివరికి అన్ని పరికల్పనలూ
\tetoken{ఉపసంహరించిన}{!!{discharged}} అయ్యాయో లేదో సులభంగా చూడవచ్చు.
\begin{prooftree}
\AxiomC{$\Discharge{!A \land !B}{1}$}
\DeduceC{$!A$}
\DischargeRule{\Intro{\lif}}{1}
\UnaryInfC{$(!A\land !B) \lif !A$}
\end{prooftree}

ఇప్పుడు $!A \land !B$ పరికల్పన నుంచి $!A$ వరకు ఉన్న మెట్లను
పూరించాలి. మనం చూసుకోవాల్సిన సంయోజకం $\land$ ఒక్కటే కనుక $\land$
తొలగింపు నియమాన్ని వాడాలి. దాంతో కింది నిరూపణ వస్తుంది:
\begin{prooftree}
\AxiomC{$\Discharge{!A \land !B}{1}$}
\RightLabel{\Elim{\land}}
\UnaryInfC{$!A$}
\DischargeRule{\Intro{\lif}}{1}
\UnaryInfC{$(!A\land !B) \lif !A$}
\end{prooftree}
ఇప్పుడు $(!A \land !B) \lif !A$కు ఒక సరైన \tetoken{వ్యుత్పత్తి}{!!{derivation}} ఉంది.
\end{ex}

\begin{ex}
ఇప్పుడు $(\lnot !A \lor !B) \lif (!A \lif !B)$కు ఒక
\tetoken{వ్యుత్పత్తి}{!!a{derivation}} ఇద్దాం.

కావలసిన నిష్కర్షను \tetoken{వ్యుత్పత్తి}{!!{derivation}} అడుగున రాయడంతో మొదలుపెడతాం.
\begin{prooftree}
\AxiomC{}
\UnaryInfC{$(\lnot !A \lor !B) \lif (!A \lif !B)$}
\end{prooftree}
ఈ నిష్కర్షను ఇవ్వగల తార్కిక నియమాన్ని కనుగొనడానికి అందులోని తార్కిక
సంయోజకాలైన $\lnot$, $\lor$, $\lif$లను చూస్తాం. ప్రస్తుతం మొదటి
$\lif$ను మాత్రమే పట్టించుకుంటాం; అది నిష్కర్షలోని \tetoken{వాక్యం}{!!{sentence}} యొక్క
\tetoken{ప్రధాన సంయోజకం}{!!{main operator}}. $\lnot$, $\lor$, రెండవ $\lif$ మరొక సంయోజకం
పరిధిలో ఉన్నాయి కనుక వాటిని తరువాత చూసుకోవచ్చు. అందువల్ల
\Intro{\lif} నియమంతో మొదలుపెడతాం. దాన్ని సరిగ్గా వర్తింపజేస్తే ఇలా
ఉండాలి:
\sourcecorrection{OLTEND-003}{మూలం సహజ నిగమన నిష్కర్షను
“అంత్య సీక్వెంట్”లోని వాక్యంగా పొరపాటుగా పేర్కొంటుంది; ఈ వ్యవస్థ
వాక్యాల వృక్షమని అధ్యాయ నిర్వచనం నిర్ణయిస్తున్నందున “నిష్కర్షలోని
వాక్యం”గా సరిచేశాం.}
\begin{prooftree}
\AxiomC{$\Discharge{\lnot !A \lor !B}{1}$}
\DeduceC{$!A \lif !B$}
\DischargeRule{\Intro{\lif}}{1}
\UnaryInfC{$(\lnot !A \lor !B) \lif (!A \lif !B)$}
\end{prooftree}
ఇక్కడి నుంచి కొనసాగడానికి రెండు మార్గాలున్నాయి. అడుగు నుంచి పైకి
పనిచేస్తూ \Intro{\lif} నియమాన్ని మరొకసారి వర్తింపజేయవచ్చు; లేదా పై
నుంచి కిందికి పనిచేస్తూ \Elim{\lor} నియమాన్ని వర్తింపజేయవచ్చు. రెండవ
మార్గాన్ని తీసుకుందాం. $\lnot !A \lor !B$ పరికల్పనను
\Elim{\lor}కు ఎడమ పూర్వాధారంగా వాడతాం. \Elim{\lor}ను చెల్లుబాటుగా
వర్తింపజేయాలంటే మిగతా రెండు పూర్వాధారాలూ నిష్కర్ష $!A \lif !B$తో
ఒకేలా ఉండాలి; అయితే వాటిలో ప్రతి ఒక్కదాన్ని $\lnot !A \lor !B$లోని
రెండు వికల్పాంగాల్లో ఒకదాన్ని పరికల్పించి వరుసగా నిగమించవచ్చు. కాబట్టి
మన \tetoken{వ్యుత్పత్తి}{!!{derivation}} ఇలా ఉంటుంది:
\begin{prooftree}
\AxiomC{$\Discharge{\lnot !A \lor !B}{1}$}
\AxiomC{$\Discharge{\lnot !A}{2}$}
\DeduceC{$!A \lif !B$}
\AxiomC{$\Discharge{!B}{2}$}
\DeduceC{$!A \lif !B$}
\DischargeRule{\Elim{\lor}}{2}
\TrinaryInfC{$!A \lif !B$}
\DischargeRule{\Intro{\lif}}{1}
\UnaryInfC{$(\lnot !A \lor !B) \lif (!A \lif !B)$}
\end{prooftree}

కుడివైపున్న రెండు శాఖల్లోనూ $!A \lif !B$ను \tetoken{నిగమనం}{!!{derive}} చేయాలి;
దానికి \Intro{\lif}ను ఉపయోగించడం ఉత్తమం.
\begin{prooftree}
\AxiomC{$\Discharge{\lnot !A \lor !B}{1}$}
\AxiomC{$\Discharge{\lnot !A}{2}, \Discharge{!A}{3}$}
\DeduceC{$!B$}
\DischargeRule{\Intro{\lif}}{3}
\UnaryInfC{$!A \lif !B$}
\AxiomC{$\Discharge{!B}{2}, \Discharge{!A}{4}$}
\DeduceC{$!B$}
\DischargeRule{\Intro{\lif}}{4}
\UnaryInfC{$!A \lif !B$}
\DischargeRule{\Elim{\lor}}{2}
\TrinaryInfC{$!A \lif !B$}
\DischargeRule{\Intro{\lif}}{1}
\UnaryInfC{$(\lnot !A \lor !B) \lif (!A \lif !B)$}
\end{prooftree}

\tetoken{వ్యుత్పత్తిలోని}{!!{derivation}} ఖాళీగా ఉన్న రెండు భాగాల కోసం, మధ్యలో $\lnot !A$,
$!A$ల నుంచి, ఎడమవైపు $!A$, $!B$ల నుంచి $!B$ యొక్క
\tetoken{వ్యుత్పత్తులు}{!!{derivation}s} కావాలి. మొదట ముందుదాన్ని తీసుకుందాం.
$\lnot !A$, $!A$లు \Elim{\lnot} యొక్క రెండు పూర్వాధారాలు:
\begin{prooftree}
\AxiomC{$\Discharge{\lnot !A}{2}$}
\AxiomC{$\Discharge{!A}{3}$}
\RightLabel{\Elim{\lnot}}
\BinaryInfC{$\lfalse$}
\DeduceC{$!B$}
\end{prooftree}
\FalseInt ఉపయోగించి $!B$ను నిష్కర్షగా పొంది శాఖను పూర్తి చేయవచ్చు.
\begin{prooftree}
\AxiomC{$\Discharge{\lnot !A \lor !B}{1}$}
\AxiomC{$\Discharge{\lnot !A}{2}$}
\AxiomC{$\Discharge{!A}{3}$}
\RightLabel{\Elim{\lnot}}
\BinaryInfC{$\lfalse$}
\RightLabel{\FalseInt}
\UnaryInfC{$!B$}
\DischargeRule{\Intro{\lif}}{3}
\UnaryInfC{$!A \lif !B$}
\AxiomC{$\Discharge{!B}{2}, \Discharge{!A}{4}$}
\DeduceC{$!B$}
\DischargeRule{\Intro{\lif}}{4}
\UnaryInfC{$!A \lif !B$}
\DischargeRule{\Elim{\lor}}{2}
\TrinaryInfC{$!A \lif !B$}
\DischargeRule{\Intro{\lif}}{1}
\UnaryInfC{$(\lnot !A \lor !B) \lif (!A \lif !B)$}
\end{prooftree}
\sourcecorrection{OLTEND-002}{మూలంలో ఈ వృక్షంలోని వైరుధ్య మెట్టుకు
పొరపాటుగా $\Intro{\lfalse}$ గుర్తు ఉంది; వెంటనే ముందున్న వివరణ,
నిర్వచిత నియమం, తర్వాతి పూర్తి వృక్షానికి అనుగుణంగా
$\Elim{\lnot}$గా సరిచేశాం.}

ఇప్పుడు కుడివైపు చివరి శాఖను చూద్దాం. \tetoken{వ్యుత్పత్తి}{!!{derivation}} నిర్వచనం
\emph{పరికల్పనలను ఉపసంహరించడానికి అనుమతిస్తుంది}, కాని వాటిని
\emph{తప్పనిసరిగా ఉపసంహరించాలని కోరదు} అని ఇక్కడ గుర్తించడం ముఖ్యం.
అంటే $!A$, $!B$ పరికల్పనల్లో ఒకదాన్ని మాత్రమే వాడి $!B$ను
నిగమించగలిగితే, మరొకదాన్ని వాడకపోవచ్చు. $!B$ నుంచి $!B$ను
\tetoken{నిగమనం}{!!{derive}} చేయడం అల్పసాధ్యం: $!B$ తనంతట తానే అలాంటి
\tetoken{ఒక వ్యుత్పత్తి}{!!a{derivation}}; ఏ నిగమనాలూ అవసరం లేదు. కాబట్టి $!A$ పరికల్పనను
తొలగించవచ్చు.
\begin{prooftree}
\AxiomC{$\Discharge{\lnot !A \lor !B}{1}$}
\AxiomC{$\Discharge{\lnot !A}{2}$}
\AxiomC{$\Discharge{!A}{3}$}
\RightLabel{\Elim{\lnot}}
\BinaryInfC{$\lfalse$}
\RightLabel{\FalseInt}
\UnaryInfC{$!B$}
\DischargeRule{\Intro{\lif}}{3}
\UnaryInfC{$!A \lif !B$}
\AxiomC{$\Discharge{!B}{2}$}
\RightLabel{\Intro{\lif}}
\UnaryInfC{$!A \lif !B$}
\DischargeRule{\Elim{\lor}}{2}
\TrinaryInfC{$!A \lif !B$}
\DischargeRule{\Intro{\lif}}{1}
\UnaryInfC{$(\lnot !A \lor !B) \lif (!A \lif !B)$}
\end{prooftree}
పూర్తయిన \tetoken{వ్యుత్పత్తిలో}{!!{derivation}} కుడివైపు చివరి \Intro{\lif} నిగమనం
వాస్తవానికి ఏ పరికల్పననూ ఉపసంహరించదని గమనించండి.
\end{ex}

\begin{ex}
ఇప్పటివరకు మనకు \FalseCl{} నియమం అవసరం కాలేదు. నియమపు నిష్కర్షకు
ఉప-\tetoken{సూత్రం}{!!{formula}} కాని ఒక పరికల్పనను ఉపసంహరించనివ్వడం దీని ప్రత్యేకత.
ఇది \FalseInt{} నియమానికి దగ్గరి సంబంధం కలది. నిజానికి \FalseInt{}
నియమం \FalseCl{} నియమపు ఒక ప్రత్యేక సందర్భం---\FalseInt{}ను మాత్రమే
అనుమతించే “అంతర్బోధాత్మక తర్కం” అనే తర్కం ఉంది. మరే పద్ధతీ పని
చేయనప్పుడు \FalseCl{} చివరి ఉపాయం. ఉదాహరణకు $!A \lor \lnot !A$ను
\tetoken{నిగమనం}{!!{derive}} చేయాలనుకుందాం. సాధారణ వ్యూహం $\Intro{\lor}$ ఉపయోగించి
$!A \lor \lnot !A$ను \tetoken{నిగమనం}{!!{derive}} చేయడానికి ప్రయత్నించడం. కానీ అలా
చేయాలంటే పరికల్పనలేమీ లేకుండా $!A$ లేదా $\lnot !A$లో ఒకదాన్ని
\tetoken{నిగమనం}{!!{derive}} చేయాలి; అది సాధ్యం కాదు. ఇప్పుడు \FalseCl{} ఆదుకుంటుంది!
\begin{prooftree}
  \AxiomC{$\Discharge{\lnot(!A \lor \lnot !A)}{1}$}
  \DeduceC{$\lfalse$}
  \DischargeRule{\FalseCl}{1}
  \UnaryInfC{$!A \lor \lnot !A$}
\end{prooftree}
ఇప్పుడు $\lnot(!A \lor \lnot !A)$ నుంచి $\lfalse$ యొక్క ఒక
\tetoken{వ్యుత్పత్తి}{!!a{derivation}} కోసం చూస్తున్నాం. $\lfalse$ అనేది $\Elim{\lnot}$
నిష్కర్ష కనుక దాన్ని ప్రయత్నించవచ్చు:
\begin{prooftree}
  \AxiomC{$\Discharge{\lnot(!A \lor \lnot !A)}{1}$}
  \DeduceC{$\lnot !A$}
  \AxiomC{$\Discharge{\lnot(!A \lor \lnot !A)}{1}$}
  \DeduceC{$!A$}
  \RightLabel{\Elim{\lnot}}
  \BinaryInfC{$\lfalse$}
  \DischargeRule{\FalseCl}{1}
  \UnaryInfC{$!A \lor \lnot !A$}
\end{prooftree}
$\lnot !A$ యొక్క ఒక \tetoken{వ్యుత్పత్తిని}{!!a{derivation}} కనుగొనే మన వ్యూహం
$\Intro{\lnot}$ను వర్తింపజేయమంటుంది:
\begin{prooftree}
  \AxiomC{$\Discharge{\lnot(!A \lor \lnot !A)}{1}, \Discharge{!A}{2}$}
  \DeduceC{$\lfalse$}
  \DischargeRule{\Intro{\lnot}}{2}
  \UnaryInfC{$\lnot !A$}
  \AxiomC{$\Discharge{\lnot(!A \lor \lnot !A)}{1}$}
  \DeduceC{$!A$}
  \RightLabel{\Elim{\lnot}}
  \BinaryInfC{$\lfalse$}
  \DischargeRule{\FalseCl}{1}
  \UnaryInfC{$!A \lor \lnot !A$}
\end{prooftree}
ఇక్కడ $\lnot(!A \lor \lnot !A)$ పరికల్పనకూ, మన కొత్త పరికల్పన
$!A$ నుంచి $\Intro{\lor}$ ద్వారా వచ్చే $!A \lor \lnot !A$కూ
$\Elim{\lnot}$ను వర్తింపజేసి $\lfalse$ను సులభంగా పొందవచ్చు:
\begin{prooftree}
  \AxiomC{$\Discharge{\lnot(!A \lor \lnot !A)}{1}$}
  \AxiomC{$\Discharge{!A}{2}$}
  \RightLabel{\Intro{\lor}}
  \UnaryInfC{$!A \lor \lnot !A$}
  \RightLabel{\Elim{\lnot}}
  \BinaryInfC{$\lfalse$}
  \DischargeRule{\Intro{\lnot}}{2}
  \UnaryInfC{$\lnot !A$}
  \AxiomC{$\Discharge{\lnot(!A \lor \lnot !A)}{1}$}
  \DeduceC{$!A$}
  \RightLabel{\Elim{\lnot}}
  \BinaryInfC{$\lfalse$}
  \DischargeRule{\FalseCl}{1}
  \UnaryInfC{$!A \lor \lnot !A$}
\end{prooftree}
కుడివైపున అదే వ్యూహాన్ని వాడతాం; అయితే $!A$ను \FalseCl ద్వారా
పొందుతాం:
\begin{prooftree}
  \AxiomC{$\Discharge{\lnot(!A \lor \lnot !A)}{1}$}
  \AxiomC{$\Discharge{!A}{2}$}
  \RightLabel{\Intro{\lor}}
  \UnaryInfC{$!A \lor \lnot !A$}
  \RightLabel{\Elim{\lnot}}
  \BinaryInfC{$\lfalse$}
  \DischargeRule{\Intro{\lnot}}{2}
  \UnaryInfC{$\lnot !A$}
  \AxiomC{$\Discharge{\lnot(!A \lor \lnot !A)}{1}$}
  \AxiomC{$\Discharge{\lnot !A}{3}$}
  \RightLabel{\Intro{\lor}}
  \UnaryInfC{$!A \lor \lnot !A$}
  \RightLabel{\Elim{\lnot}}
  \BinaryInfC{$\lfalse$}
  \DischargeRule{\FalseCl}{3}
  \UnaryInfC{$!A$}
  \RightLabel{\Elim{\lnot}}
  \BinaryInfC{$\lfalse$}
  \DischargeRule{\FalseCl}{1}
  \UnaryInfC{$!A \lor \lnot !A$}
\end{prooftree}
\end{ex}

\begin{prob}
కిందివాటిని చూపే \tetoken{వ్యుత్పత్తులు}{!!{derivation}s} ఇవ్వండి:
\begin{enumerate}
\item $!A \land (!B \land !C) \Proves (!A \land !B) \land !C$.
\item $!A \lor (!B \lor !C) \Proves (!A \lor !B) \lor !C$.
\item $!A \lif (!B \lif !C) \Proves !B \lif (!A \lif !C)$.
\item $!A \Proves \lnot\lnot !A$.
\end{enumerate}
\end{prob}

\begin{prob}
కిందివాటిని చూపే \tetoken{వ్యుత్పత్తులు}{!!{derivation}s} ఇవ్వండి:
\begin{enumerate}
\item $(!A \lor !B) \lif !C \Proves !A \lif !C$.
\item $(!A \lif !C) \land (!B \lif !C) \Proves (!A \lor !B) \lif !C$.
\item $\Proves \lnot(!A \land \lnot !A)$.
\item $!B \lif !A \Proves \lnot !A \lif \lnot !B$.
\item $\Proves (!A \lif \lnot !A) \lif \lnot !A$.
\item $\Proves \lnot(!A \lif !B) \lif \lnot !B$.
\item $!A \lif !C \Proves \lnot (!A \land \lnot !C)$.
\item $!A \land \lnot !C \Proves \lnot (!A \lif !C)$.
\item $!A \lor !B, \lnot !B \Proves !A$.
\item $\lnot !A \lor \lnot !B \Proves \lnot(!A \land !B)$.
\item $\Proves (\lnot !A \land \lnot !B) \lif\lnot(!A \lor !B)$.
\item $\Proves \lnot(!A \lor !B) \lif (\lnot !A \land \lnot !B)$.
\end{enumerate}
\end{prob}

\begin{prob}
కిందివాటిని చూపే \tetoken{వ్యుత్పత్తులు}{!!{derivation}s} ఇవ్వండి:
\begin{enumerate}
\item $\lnot(!A \lif !B) \Proves !A$.
\item $\lnot(!A \land !B) \Proves \lnot !A \lor \lnot !B$.
\item $!A \lif !B \Proves \lnot !A \lor !B$.
\item $\Proves \lnot \lnot !A \lif !A$.
\item $!A \lif !B, \lnot !A \lif !B \Proves !B$.
\item $(!A \land !B) \lif !C \Proves (!A \lif !C) \lor (!B \lif !C)$.
\item $(!A \lif !B) \lif !A \Proves !A$.
\item $\Proves (!A \lif !B) \lor (!B \lif !C)$.
\end{enumerate}
(వీటన్నింటికీ $\FalseCl$~నియమం అవసరం.)
\end{prob}

\end{document}
